% --- ARCHIVO: Cap8_Uso_IA/uso_ia.tex ---

\chapter{Discusión sobre el uso de Inteligencia Artificial en el TFG}
\label{cap:uso_ia}

La elaboración de este Trabajo de Fin de Grado se ha desarrollado en un contexto
metodológico atípico: la concurrencia simultánea, durante el bienio 2024--2026,
de una crisis sistémica en el sistema eléctrico ibérico y de una transformación
acelerada en las herramientas de procesamiento de lenguaje natural. La
disponibilidad de \textit{Large Language Models} (LLMs) capaces de procesar
documentos técnicos de cientos de páginas en pocos minutos coincidió con la
publicación, en el plazo de un año, de cuatro informes periciales primarios y
una decena de análisis académicos sobre el mismo evento. El presente capítulo
documenta de forma transparente cómo se han integrado estas herramientas en el
flujo de trabajo, qué tareas concretas han desempeñado, qué errores
sistemáticos se han detectado y corregido, y qué reflexión epistemológica se
desprende de todo ello para la práctica investigadora en ingeniería.

\section{Aplicación metodológica: Uso de IA en la estructuración, síntesis de
informes masivos y extracción de datos}
\label{sec:aplicacion_ia}

El cuerpo documental de partida abarca expedientes técnicos, normativas de
ENTSO-E, informes oficiales de extensión superior a las 200~páginas,
simulaciones PSS/E reproducidas por el IIT-ICAI, actas de la CNMC, registros
oscilográficos publicados por el NREL y un volumen amplio de hemeroteca
especializada. Su lectura lineal sin mediación informática resultaba
incompatible con los plazos académicos del TFG. La integración de modelos
generativos de gran escala se ha articulado, en consecuencia, en torno a cinco
funciones operativas claramente delimitadas, ninguna de las cuales involucra
inferencia causal autónoma sobre la física del sistema.

La primera función ha sido la \textbf{reconciliación cronológica} entre
informes con sistemas de sellado temporal distintos. El informe del Consejo de
Seguridad Nacional, el de Red Eléctrica de España, el del IIT-ICAI y el factual
de ENTSO-E describen los mismos hechos físicos del 28 de abril de 2025 con
distintos niveles de granularidad temporal y, ocasionalmente, con desfases de
hasta 200~ms entre versiones. La asistencia automática ha permitido construir
una tabla maestra unificada de \textit{timestamps} y referencias cruzadas que
ha servido de andamiaje para los capítulos~\ref{cap:analisis_incidente}
(p.~\pageref{cap:analisis_incidente}) y \ref{cap:reaccion_reposicion}
(p.~\pageref{cap:reaccion_reposicion}).

La segunda función ha sido el \textbf{mapeo sistemático de divergencias} entre
narrativas. Cada informe pericial enmarca el incidente en un perímetro
analítico distinto ---cumplimiento normativo, estabilidad capacitiva,
oscilaciones inter-área, diseño de mercado--- y los puntos de fricción no son
siempre evidentes a primera lectura. La extracción asistida de pasajes
homólogos ha permitido localizar las afirmaciones técnicamente incompatibles
entre fuentes, que constituyen la base del capítulo~\ref{cap:analisis_informes}
(p.~\pageref{cap:analisis_informes}) y de la síntesis comparativa de la
sección~\ref{sec:sintesis_agentes} (p.~\pageref{sec:sintesis_agentes}).

La tercera función ha sido la \textbf{tabulación cuantitativa} de magnitudes
distribuidas a lo largo de los informes: valores zonales de inercia,
penetración renovable instantánea, intercambios transfronterizos, cifras de
deslastre por subfrecuencia (UFLS), aportaciones capacitivas atribuidas a la
maniobra de mallado. Estas cifras dispersas, una vez consolidadas, han
permitido construir las gráficas y tablas que articulan visualmente la
narrativa técnica del trabajo.

La cuarta función ha sido el \textbf{soporte redaccional} en español técnico:
sugerencia de variantes léxicas, depuración de sinonimias, control de
terminología y reformulación de pasajes de prosa académica densa. Este apoyo
no ha alterado en ningún caso la estructura argumental ni las decisiones de
contenido, que son responsabilidad exclusiva del autor.

La quinta función, finalmente, ha sido el \textbf{andamiaje de elementos
gráficos}: generación de borradores en TikZ y \texttt{pgfplots} a partir de
descripciones textuales, posteriormente revisados, depurados y validados
manualmente para asegurar consistencia visual con el resto del documento y
fidelidad a los datos representados.

\begin{figure}[H]
    \centering
    \resizebox{\textwidth}{!}{%
    \begin{tikzpicture}[
        node distance=0.7cm and 0.9cm,
        every node/.style={font=\footnotesize},
        block/.style={rectangle, rounded corners=3pt, draw=black!70, thick,
                      fill=gray!10, align=center, minimum width=2.6cm,
                      minimum height=1.0cm, inner sep=3pt},
        ai/.style={rectangle, rounded corners=3pt, draw=blue!60!black, thick,
                   fill=blue!8, align=center, minimum width=2.8cm,
                   minimum height=1.0cm, inner sep=3pt},
        filter/.style={rectangle, rounded corners=3pt, draw=red!70!black,
                       thick, fill=red!8, align=center, minimum width=3.2cm,
                       minimum height=1.0cm, inner sep=3pt},
        outbox/.style={rectangle, rounded corners=3pt, draw=green!50!black,
                       thick, fill=green!10, align=center, minimum width=2.6cm,
                       minimum height=1.0cm, inner sep=3pt},
        arr/.style={->, thick, gray!70},
        loop/.style={->, thick, red!70!black, dashed}
    ]
        \node[block] (src1) {Informe CSN\\(\textit{Gobierno})};
        \node[block, below=of src1] (src2) {Informe REE};
        \node[block, below=of src2] (src3) {Informe IIT-ICAI\\AELEC};
        \node[block, below=of src3] (src4) {Informe ENTSO-E\\NREL, MIT};

        \node[ai, right=of src2, yshift=-0.55cm] (ia)
            {\textbf{Asistencia LLM}\\
            {\scriptsize\itshape Extracci\'on -- S\'intesis}\\
            {\scriptsize\itshape Tabulaci\'on -- Cronolog\'ia}};

        \node[filter, right=of ia] (phys)
            {\textbf{Filtro de validaci\'on f\'isica}\\
            {\scriptsize Leyes de Kirchhoff -- curvas $Q$-$V$}\\
            {\scriptsize Ecuaci\'on oscilaci\'on del rotor}\\
            {\scriptsize Cotejo con fuente primaria}};

        \node[outbox, right=of phys] (outnode)
            {\textbf{Resultado}\\\textbf{validado}\\
            {\scriptsize\itshape (cuerpo del TFG)}};

        \draw[arr] (src1.east) -- (ia.west);
        \draw[arr] (src2.east) -- (ia.west);
        \draw[arr] (src3.east) -- (ia.west);
        \draw[arr] (src4.east) -- (ia.west);

        \draw[arr] (ia) -- (phys);
        \draw[arr] (phys) -- (outnode);

        \draw[loop] (phys.south) -- ++(0,-0.9) -| (ia.south);
        \node[font=\scriptsize\itshape, text=red!70!black, align=center,
              below=1.1cm of ia.south, xshift=2.2cm]
             {Reformulaci\'on de \textit{prompt}\\ante alucinaci\'on detectada};

    \end{tikzpicture}%
    }
    \caption[Flujo de trabajo de la asistencia por IA con bucle de validaci\'on f\'isica.]{Flujo de trabajo metodol\'ogico empleado en el TFG. Las fuentes primarias alimentan a la asistencia LLM, cuyo resultado pasa obligatoriamente por un filtro de validaci\'on f\'isica antes de incorporarse al cuerpo del trabajo. Cuando el filtro detecta una inferencia incompatible con la mec\'anica del sistema (alucinaci\'on), el \textit{prompt} se reformula con restricciones sem\'anticas adicionales y el ciclo se reinicia. Fuente: elaboraci\'on propia}
    \label{fig:workflow_ia}
\end{figure}

El elemento metodológicamente decisivo de este flujo no es la asistencia LLM
en sí misma, sino el \textbf{filtro de validación física} que se interpone
entre cualquier salida del modelo y su incorporación al cuerpo del trabajo.
La sección siguiente documenta los casos en los que ese filtro detuvo
inferencias erróneas que, de haberse incorporado al texto sin revisión,
habrían comprometido la integridad técnica del análisis.

\section{Validación y Verificación: Proceso de contraste riguroso de los
resultados de la IA con la física del sistema y la normativa oficial}
\label{sec:validacion_ia}

La utilización de modelos generativos en un dominio tan denso en relaciones
causales como la dinámica de sistemas de potencia revela, de forma sistemática,
un comportamiento que los sistemas eléctricos no toleran: la inferencia por
analogía estadística. Los LLM, entrenados sobre grandes corpus generales,
tienden a aplicar patrones causales \textit{típicos} ---los más frecuentes en
la literatura de su corpus de entrenamiento--- a casos cuya estructura física
es \textit{atípica}. El 28 de abril de 2025 es, por su naturaleza dual
(estabilidad de tensión en lugar de balance de frecuencia), un caso atípico
respecto a la mayoría de los apagones documentados en la bibliografía técnica
del siglo~XX, lo que ha hecho especialmente productiva la observación de los
puntos en los que la IA «por defecto» se desviaba del fenómeno real.

La tabla~\ref{tab:alucinaciones_ia} sintetiza los casos más relevantes de
inferencia errónea detectados durante el desarrollo del trabajo, junto con la
corrección física aplicada y el mecanismo de \textit{prompt} restringido
empleado para depurarla.

\begin{table}[H]
    \centering
    \caption[Casos representativos de alucinación detectados y corrección
    física aplicada.]{Casos representativos de inferencia errónea detectada en
    los modelos generativos durante el desarrollo del trabajo, junto con la
    corrección física aplicada y la estrategia de \textit{prompt} restringido
    utilizada. Fuente: elaboración propia.}
    \label{tab:alucinaciones_ia}
    \vspace{0.2cm}
    \footnotesize
    \renewcommand{\arraystretch}{1.25}
    \begin{tabularx}{\textwidth}{@{}>{\bfseries\raggedright\arraybackslash}p{2.6cm}
                                  >{\raggedright\arraybackslash}X
                                  >{\raggedright\arraybackslash}X
                                  >{\raggedright\arraybackslash}X@{}}
        \toprule
        \textnormal{\textbf{Fenómeno}} &
        \textbf{Inferencia errónea \textit{por defecto} del LLM} &
        \textbf{Corrección física aplicada} &
        \textbf{Estrategia de \textit{prompt} correctora} \\
        \midrule

        Paradoja del UFLS &
        El deslastre de cargas por subfrecuencia «salvó» áreas del sistema al
        recuperar el balance de potencia activa. &
        El UFLS es ciego al voltaje: al desconectar carga retiró sumideros de
        reactiva inductiva en pleno transitorio capacitivo, agravando la
        sobretensión. &
        Restricción explícita: razonar exclusivamente en el plano $Q$-$V$,
        ignorando la lógica frecuencia-potencia activa habitual en apagones
        clásicos. \\

        \addlinespace[2pt]
        \textit{Tap-Lag} &
        Los OLTC respondieron correctamente al transitorio elevando o
        rebajando la tensión secundaria de forma proporcional. &
        Los retardos intencionales de los OLTC ---diseñados para evitar
        \textit{hunting}--- los dejaron desfasados frente a un transitorio de
        decenas de milisegundos, multiplicando la tensión hacia la red de
        220~kV y 132~kV. &
        Inyección de la constante de tiempo real del conjunto electromecánico
        OLTC y prohibición de asumir respuesta cuasi-instantánea del cambiador
        de tomas. \\

        \addlinespace[2pt]
        Maniobra de mallado y aporte capacitivo &
        El mallado redujo la impedancia del sistema y, por tanto, aumentó la
        estabilidad de tensión. &
        En condiciones de baja carga, el mallado activó la admitancia
        transversal de las líneas de 400~kV (efecto Ferranti), inyectando los
        1.050~MVAr capacitivos que precipitaron la sobretensión. &
        Restricción al régimen «línea descargada» y exigencia de calcular el
        balance de reactiva con la admitancia capacitiva $Y_t$ en el modelo
        $\pi$ de la línea. \\

        \addlinespace[2pt]
        Oscilación de 0,6~Hz &
        Modo electromecánico inter-área clásico análogo a los modos de 0,2~Hz
        del sistema síncrono europeo. &
        Modo de origen controvertido entre informes (forzado por planta
        generadora según ICAI; modo natural del sistema según REE), que el
        TFG presenta de forma comparada sin resolución editorial. &
        Veto a la elección unilateral entre ambas hipótesis y obligación de
        presentar las dos lecturas con atribución a su fuente. \\

        \bottomrule
    \end{tabularx}
\end{table}

El caso de la \textbf{Paradoja del UFLS} ilustra con especial claridad el
patrón de error. Ante la consulta sobre el papel del deslastre de carga en el
desarrollo del incidente, los modelos consultados ofrecieron sistemáticamente
una interpretación heredada de los apagones clásicos, en los que el UFLS
constituye la última línea de defensa y consigue habitualmente preservar
porciones de la red al recuperar el balance frecuencia-potencia activa. Esta
inferencia es físicamente correcta para el universo de eventos sobre los que
el modelo ha sido entrenado, pero resulta inválida para el 28A: en un colapso
dominado por el plano $Q$-$V$, la activación del UFLS retiró simultáneamente
los sumideros de potencia reactiva ($Q$) inductiva que aportaban las cargas
desconectadas, lo que agravó la sobretensión en lugar de mitigarla. La
corrección requirió formular el \textit{prompt} con una restricción explícita
sobre el plano de análisis ---razonar en términos de tensión y reactiva,
ignorando el reflejo automático hacia el dominio de frecuencia--- y verificar
manualmente que el razonamiento resultante fuese consistente con la
oscilografía publicada por el NREL en la franja temporal 12:33:19--12:33:27
\cite{nrel2025, icai2025}.

El caso del \textbf{Tap-Lag} reveló una limitación distinta pero del mismo
género. La descripción cualitativa que los modelos producen del cambiador de
tomas en carga (OLTC) es correcta en régimen estacionario, pero asume por
defecto una respuesta proporcional e instantánea que el conjunto
electromecánico real no posee. El retardo intencional introducido para evitar
oscilación mecánica del cambiador (\textit{hunting}) se tradujo el 28A en una
relación de transformación desfasada respecto al transitorio eléctrico, lo que
multiplicó la sobretensión hacia las redes de 220~kV y 132~kV
---fenómeno invisible para el SCADA del operador, anclado en las medidas del
primario de 400~kV---. Toda la cadena causal del Tap-Lag exigió describir al
modelo, por adelantado, las constantes de tiempo del OLTC y prohibirle asumir
respuesta cuasi-instantánea \cite{ree2025, gobierno2025}.

El caso del \textbf{aporte capacitivo del mallado} planteó la verificación
cuantitativa más exigente. La inyección estimada de 1.050~MVAr capacitivos
asociada a la maniobra de mallado debía corresponderse con el efecto de la
admitancia transversal ($Y_t$) de las líneas de 400~kV operando en régimen
descargado, fenómeno conocido como efecto Ferranti. La validación consistió
en reproducir el balance de reactiva sobre el modelo en $\pi$ de la línea,
contrastando la magnitud orden de grandeza con la oscilografía y con el
diagnóstico independiente del informe IIT-ICAI \cite{icai2025}.

Por último, el caso de la \textbf{oscilación de 0,6~Hz a las 12:03~CEST}
ilustra una limitación específica del uso de IA en contextos forenses con
divergencia institucional. Los modelos tienden a producir una respuesta
\textit{integrada} ---la lectura «más probable» promediada sobre las
fuentes--- cuando el rigor pericial exige preservar la divergencia: el origen
del modo (oscilación forzada por una planta concreta, según el IIT-ICAI;
modo natural del sistema agravado por el estado operativo, según REE) es
precisamente el dato que el TFG debe preservar como discrepancia, no resolver
editorialmente. La estrategia correctora exigió vetar al modelo cualquier
elección unilateral entre las dos hipótesis y obligarle a producir versiones
duales con atribución explícita \cite{icai2025, ree2025}.

El conjunto de estos casos permite formular una observación de alcance general
sobre el empleo de LLM en investigación de ingeniería: la utilidad de la
asistencia automática es máxima en tareas de \textit{procesamiento}
---extracción, síntesis, estructuración, tabulación, redacción asistida---
y mínima, hasta el punto de ser contraproducente, en tareas de
\textit{inferencia causal} sobre fenómenos físicos con baja representación en
los corpus de entrenamiento. La dimensión metodológicamente más exigente del
trabajo no ha consistido, por tanto, en aprovechar las capacidades del modelo,
sino en \textbf{reconocer activamente sus límites} y construir un protocolo de
validación capaz de detenerse en el momento exacto en que la maquinaria
estadística del LLM se desviaba de la física del caso.

\section{Reflexión crítica sobre la IA en la ingeniería de investigación}
\label{sec:reflexion_ia}

La experiencia documentada en las dos secciones anteriores invita a una
reflexión epistemológica más amplia sobre el papel que estas herramientas
están llamadas a desempeñar en la investigación en ingeniería. La conclusión
que se desprende del trabajo desarrollado no es triunfalista ni
catastrofista, sino estructural: los modelos generativos son potentes
\textit{aceleradores} de los procesos cognitivos del investigador, pero
operan sobre una lógica probabilística que no coincide con la lógica
determinista que rige la física de los sistemas materiales. La integridad
del análisis depende, por tanto, de mantener visible y operativa esa
distinción a lo largo de todo el trabajo.

La primera implicación, de carácter \textbf{metodológico}, es que el uso de
IA no exime al investigador de dominar el contenido técnico de su objeto de
estudio; al contrario, lo exige con mayor severidad. Solo un investigador con
competencia técnica suficiente para identificar una inferencia físicamente
imposible puede emplear con seguridad un modelo capaz de producirla con
fluidez retórica. La asimetría entre la verosimilitud lingüística del
\textit{output} y su corrección material constituye el principal riesgo
asociado al uso de estas herramientas en dominios técnicos exigentes, y solo
puede neutralizarse mediante un dominio del contenido que la herramienta no
puede aportar.

La segunda implicación, de carácter \textbf{ético-académico}, se refiere a
la atribución de autoría intelectual. La asistencia automática se ha
empleado en este trabajo como herramienta de procesamiento, no como
sustituto del razonamiento. Toda decisión de contenido ---el encuadre del
problema, la jerarquía de fuentes, la formulación de la tesis causal sobre
el colapso del 28A, la valoración comparada de las narrativas
institucionales y las conclusiones del último capítulo--- es responsabilidad
exclusiva del autor. La IA no firma TFGs ni se sienta ante un tribunal: el
ingeniero rubrica con su nombre la causalidad material de los hechos
expuestos, y esa rúbrica es intransferible.

La tercera implicación, de carácter \textbf{disciplinar}, abre una línea
prospectiva. Los sistemas eléctricos de la próxima década ---dominados por
electrónica de potencia, con dinámica relevante en la escala sub-cíclica y
con interacciones tensión--reactiva que el operador convencional no captura
en tiempo real--- van a requerir herramientas de soporte a la decisión que
combinen la velocidad de procesamiento de los modelos algorítmicos con la
robustez física de los simuladores tipo PSS/E. La integración no es trivial:
los simuladores resuelven ecuaciones; los LLM aproximan distribuciones de
probabilidad. La arquitectura híbrida ---LLM como interfaz semántica y de
síntesis, simuladores físicos como motor de inferencia causal--- es,
plausiblemente, la dirección hacia la que evolucionarán las herramientas de
ingeniería de sistemas de potencia. Su desarrollo, sin embargo, requerirá
exactamente la misma vigilancia epistemológica que ha estructurado el flujo
de trabajo del presente TFG: tener siempre presente qué hace cada
herramienta y qué se le pide explícitamente que no haga.

El registro detallado del flujo de \textit{prompts} restrictivos empleados,
junto con ejemplos representativos de las correcciones aplicadas, se
documenta en el anexo~\ref{app:prompts}
(p.~\pageref{app:prompts}), siguiendo las recomendaciones de buenas
prácticas en transparencia metodológica que la literatura académica reciente
viene consolidando para la integración de modelos generativos en
investigación de grado y posgrado.